\section{Introduction}

\subsection{Motivation}

Transformer self-attention has quadratic complexity in sequence length, 
motivating a large body of work on efficient alternatives. **State-space 
models\textbf{ (SSMs), particularly }Mamba** (Gu \& Dao, 2023), have emerged as a 
compelling family: they achieve $\mathcal{O}(L)$ complexity via selective 
state updates with input-dependent dynamics.

However, replacing a pretrained transformer's attention with an SSM is 
non-trivial:

\begin{enumerate}
\item \textbf{Pretrained weights:} The SSM's input/output projections have no
  correspondence to the pretrained $W_Q, W_K, W_V, W_O$.
\item \textbf{Residual flow:} SSMs process the input in a fundamentally different way,
  changing the residual stream statistics.
\item \textbf{Training cost:} Training a full SSM from scratch is expensive.
\item \textbf{Causal vs bidirectional:} SSMs are naturally causal; encoders need
  bidirectional processing.

\end{enumerate}
\textbf{Mamba-Inspired Gated State-Space Attention} addresses these challenges:

\begin{itemize}
\item It \textbf{reuses pretrained weights} by projecting Q, K, V into a two-branch
  gating structure.
\item It applies a \textbf{learned exponential decay} to emulate selective state-space
  dynamics without a true SSM scan.
\item It adds a \textbf{preservation gate} that blends the new output with the residual
  input, stabilizing the initial forward pass.
\item It supports both causal (decoders) and bidirectional (encoders) modes.

\end{itemize}
\subsection{Problem Statement}

Given a pretrained transformer $\mathcal{M}$ with attention projections 
$\{W_Q, W_K, W_V, W_O\}$, we seek a replacement mechanism $\mathcal{A}'$ such 
that:

\[
\mathcal{A}' \approx \mathcal{A}, \qquad
\mathcal{A}' \text{ has SSM-inspired structure}
\tag{1.1}
\]

subject to:

\begin{enumerate}
\item \textbf{Weight reuse:} Must copy pretrained weights via \texttt{copy\_weights}.
\item \textbf{Multi-architecture:} DistilBERT, RoBERTa, GPT-2, Qwen.
\item \textbf{Causal and bidirectional:} Support both modes.
\item \textbf{Recovery capacity:} LoRA adapters for fine-tuning.
\item \textbf{Preservation:} Blend with residual input to stabilize.

\end{enumerate}
\subsection{Our Approach}

We propose a gated state-space-inspired attention with:

\subsubsection{Two-Branch In-Projection}

The input projection is doubled: $W_\text{in} \in \mathbb{R}^{d \times 2d}$. 
The output splits into:

\begin{itemize}
\item $p_1$: SSM input representation (projected from $W_V$ or $W_Q$)
\item $p_2$: gating representation (projected from $W_Q$ or $W_K$)

\end{itemize}
\subsubsection{Learned Exponential Decay}

For each dimension $i$, we learn a decay $\gamma_i \in (0, 1)$ via:

\[
\gamma_i = \sigma(-\exp(A_{\log, i}))
\tag{1.2}
\]

The pairwise weight between positions $j$ and $i$ ($j \leq i$):

\[
w_{ij}^{(d)} = \frac{\exp(\text{dist}_{ij} \cdot \log \gamma_d)}{\sum_{j'} \exp(\text{dist}_{ij'} \cdot \log \gamma_d)}
\tag{1.3}
\]

where $\text{dist}_{ij} = \max(0, j - i)$ for causal mode.

\subsubsection{Gated Fusion}

\[
\text{Out} = \text{SiLU}(p_1^\text{decayed}) \odot p_2
\tag{1.4}
\]

\subsubsection{Preservation Gate}

\[
\text{Final} = \sigma(g) \cdot \text{Out} + (1 - \sigma(g)) \cdot X
\tag{1.5}
\]

with $g$ initialized to 2.0 (so $\sigma(g) \approx 0.88$).

\subsection{Contributions}

\begin{itemize}
\item \textbf{A gated SSM-inspired attention.} Reuses pretrained Q/K/V/O weights in a
  two-branch gated state-space formulation with learned exponential decay.
\item \textbf{A preservation gate.} Blends the SSM output with the residual input,
  stabilizing the initial forward pass.
\item \textbf{Multi-architecture support.} Handles DistilBERT, RoBERTa, GPT-2, and
  Qwen.
\item \textbf{Causal and bidirectional modes.} Learned decay for encoders uses both
  forward and backward directions.
\item \textbf{Honest complexity reporting.} We acknowledge that the implementation is
  $\mathcal{O}(L^2 D)$, \textbf{not} $\mathcal{O}(L)$, despite the SSM inspiration.

\end{itemize}
\subsection{Paper Organization}

Section 2 reviews related work on SSMs and efficient attention. Section 3 
provides background on state-space models, Mamba, and attention. Section 4 
presents our method. Section 5 describes integration. Section 6 details the 
setup. Section 7 reports results. Section 8 analyzes the SSM-inspired 
behavior. Section 9 discusses limitations and trade-offs. Section 10 lists 
limitations. Section 11 concludes.



\subsection{Notation}

\begin{lstlisting}

---
\end{lstlisting}